%-----------------------------------------------------------------------
\subsection{Syntactic Ambiguity (Linguistics)}
\label{sec:instance-syntax}
%-----------------------------------------------------------------------

Syntactic ambiguity is pervasive in natural language
\citep{chomsky1957,jurafsky2024}.  The classic example is prepositional phrase
(PP) attachment: the sentence fragment ``V NP PP'' admits two parse trees:
\begin{itemize}
  \item \textbf{NP-attachment} (low attachment): $(\text{V}\;(\text{NP}\;\text{PP}))$
    --- the PP modifies the NP;
  \item \textbf{VP-attachment} (high attachment): $((\text{V}\;\text{NP})\;\text{PP})$
    --- the PP modifies the verb.
\end{itemize}
For instance, ``saw the man with the telescope'' can mean either ``saw
[the man who had the telescope]'' or ``[saw the man] [using the telescope].''
Both bracketings yield the same surface token sequence $[\text{V}, \text{NP},
\text{PP}]$ but assign different constituent structures.

\paragraph{The explanation system.}
Define $\mathcal{S}_{\mathrm{syntax}} = (\Params, \sH, Y, \mathsf{obs},
\mathsf{exp}, \incomp)$ as follows:
\begin{itemize}
  \item $\Params = \sH = \texttt{Bracketing}$ (parse trees);
  \item $Y = \texttt{List Token}$ (the surface token sequence);
  \item $\mathsf{obs} = \texttt{yield}$ (the linearization of the parse tree);
  \item $\mathsf{exp} = \mathrm{id}$ (the parse tree is the ``explanation''
    of the sentence);
  \item $\incomp(b_1, b_2) \iff b_1 \ne b_2$.
\end{itemize}

\paragraph{The Rashomon property.}
The witnesses are \texttt{leftAttach} and \texttt{rightAttach}, which satisfy
$\texttt{yield}(\texttt{leftAttach}) = [\text{V}, \text{NP}, \text{PP}] =
\texttt{yield}(\texttt{rightAttach})$ and
$\texttt{leftAttach} \ne \texttt{rightAttach}$.
The Rashomon property is constructively witnessed with zero axioms.

\paragraph{The impossibility.}
By Theorem~\ref{thm:universal}, no parser can simultaneously be faithful
(return the actual parse tree), stable (assign the same parse to all token
sequences with the same yield), and decisive (distinguish every pair of
distinct parse trees).  In computational linguistics: when a sentence is
genuinely ambiguous, no deterministic parser can correctly report the syntactic
structure without either being unfaithful to one reading or unstable across
equivalent inputs.

\paragraph{Empirical illustration.}
Four syntactic parsers (spaCy small/medium/large, Stanza) applied to 50
ambiguous and 50 unambiguous sentences show mean inter-parser agreement (UAS) of
0.821 for ambiguous sentences versus 0.894 for unambiguous (Wilcoxon
$p = 1.6 \times 10^{-3}$, Cohen's $d = 0.60$).
\emph{Caveats:} Three of the four parsers (spaCy variants) share training data
and architecture, inflating within-family agreement.  The ambiguous set includes
mixed ambiguity types (PP-attachment, coordination scope, headline ambiguity),
not exclusively PP-attachment.  Parsers tokenize independently, introducing a
confound.  These limitations weaken the empirical illustration but do not affect
the Lean-verified impossibility, which depends only on the constructive witness.

\paragraph{Resolution: packed parse forests.}
The standard resolution in NLP is to report a \emph{packed parse forest}
\citep{jurafsky2024}: a compact representation of all valid parse trees for
the input.  This is the $G$-invariant aggregation over the symmetry group of
syntactic rearrangements within the same yield fiber.  Chart parsers (CYK,
Earley) naturally produce packed forests.  The resolution sacrifices
decisiveness (no single tree is selected) to achieve stability and
faithfulness.

\paragraph{Lean formalization.}
The Lean~4 types are \texttt{Token} (an inductive with constructors \texttt{V},
\texttt{NP}, \texttt{PP} representing phrasal categories, not surface tokens),
\texttt{Bracketing} (with constructors \texttt{leftAttach} and
\texttt{rightAttach}, corresponding to NP-attachment and VP-attachment), and
\texttt{syntaxSystem : ExplanationSystem Bracketing Bracketing (List Token)}.
The impossibility is \texttt{syntactic\_ambiguity\_impossibility} in
\texttt{SyntacticAmbiguity.lean}, derived with zero axioms.
